% !TeX root = ../../Book4.tex
% 纲要标题：## 附录 B　正合范畴、K 理论与更高同调结构
% 纲要位置：工作记录/计划纲要.md，第 3862 行。
% 纲要细目（写作范围）：
% > 正合结构、低阶 K 群与完美复形的可加性给出完整证明和算例；一般正合导出范畴、Q 构造的比较与局部化定理给出证明概要，模型与稳定无穷范畴补充构造及论证概要。

\chapter{正合范畴、K 理论与更高同调结构}
\label{b4:app:B}

\begin{chapterabstract}
本附录从正合结构出发建立代数 \(K_0\) 与矩阵 \(K_1\) 的基本计算，比较投射模和完美复形的 Euler 类，并介绍高阶 \(K\) 群与同伦模型的定义范围。
\end{chapterabstract}

\begin{learninggoals}
\begin{itemize}
\item 区分加性范畴、指定正合结构与交换范畴，并核验分裂和扩张封闭的例子。
\item 判断正合子范畴中的零调条件，说明分裂正合复形为何可缩。
\item 利用 \(K_0\) 的泛性质、Euler 类和三角关系进行计算。
\item 用初等矩阵证明稳定交换子公式，并求 \(K_1(k)\) 与 \(K_1(\mathbb Z)\)。
\item 写出 \(Q\) 构造的态射、复合和高阶定义，说明低阶比较与局部化长正合列的构造。
\item 在复形上核验模型结构，并辨认普通同伦态射所遗失的高阶映射信息。
\end{itemize}
\end{learninggoals}

有限生成投射模对直和与直和因子都很稳定，却未必对余核稳定：\(\mathbb Z\xrightarrow{2}\mathbb Z\) 的两端都是自由模，余核 \(\mathbb Z/2\mathbb Z\) 却不是投射模。因此若只在有限生成投射模之间工作，不能照搬模范畴的全部短正合列。我们可以先指定哪些序列允许充当「正合序列」，再为每条获准的 \(0\rightarrow A\rightarrow B\rightarrow C\rightarrow0\) 规定 \([B]=[A]+[C]\)，记录其中的加法关系。这个问题把正合结构与 \(K_0\) 联系起来；继续比较矩阵自同构与更高同伦信息，就会走向 \(K_1\) 和高阶 \(K\) 群。

\ProseParagraphBreak
所需基础为加性与交换范畴、短正合列、投射模、复形与链同伦；讨论完美复形时还使用本卷已建立的三角范畴和有界投射模型。

\ProseParagraphBreak
\chapref{b4:ch:04}的正合性问题在\secref{b4:sec:B01}发展为正合结构；\secref{b4:sec:B02}研究投射模与完美复形的可加不变量，形成完整的 \(K_0\) 单元。\subsecref{b4:subsec:B03:01}接着以矩阵计算 \(K_1\)。\secref{b4:sec:B03}余下的 \(Q\) 构造及\secref{b4:sec:B04}需要神经、几何实现和同伦群的概念。

\ProseParagraphBreak
主要参考为 Weibel《The K-book》\cite[Chapter II, Section 7; Chapter III, Section 1; Chapter IV, Sections 6--7; Chapter V, Theorem 5.1 and Definition 10.4]{Weibel2013KBook}，正合范畴的系统讨论见 Bühler《Exact Categories》\cite[Sections 2 and 10]{Buehler2010ExactCategories}。稳定无穷范畴的阅读范围限于 Blumberg、Gepner、Tabuada\cite[Section 2.2, pp. 747--750]{BlumbergGepnerTabuada2013UniversalKTheory}中相关定义与模型联系。

% !TeX root = ../../Book4.tex
% 纲要标题：### B.1 正合范畴与导出范畴
% 纲要位置：工作记录/计划纲要.md，第 3569 行。
% 纲要细目（写作范围）：
% #### B.1.1 正合结构的公理与分裂情形
% #### B.1.2 扩张封闭子范畴与可容许正合列
% #### B.1.3 零调复形与有界导出范畴

\section{正合范畴与导出范畴}
\label{b4:sec:B01}

有限生成投射模之间有直和、核的某些特例和分裂短正合列，却通常没有全部余核。若要求它们组成交换范畴，就必须扩大对象范围。正合范畴保留原来的加性范畴，只指定其中哪些列按短正合列使用。

\subsection{正合结构的公理与分裂情形}
\label{b4:subsec:B01:01}

\begin{definition}\label{b4:def:B-exact-category}
设 \(\mathcal E\) 为加性范畴。选取一族态射对 \(A\xrightarrow{i}B\xrightarrow{p}C\)，其中 \(i\) 是 \(p\) 的核、\(p\) 是 \(i\) 的余核，并要求：
\begin{enumerate}
\item 该族在列同构下封闭，包含全部分裂列 \(A\rightarrow A\oplus C\rightarrow C\)。
\item 出现在这些列中的第一个态射的复合仍是这种态射；第二个态射的复合也如此。
\item 第一个态射沿任意态射作推出存在，推出后仍为第一个态射；第二个态射沿任意态射作拉回存在，拉回后仍为第二个态射。
\end{enumerate}
这样的指定称为一个\term[zhenghejiegou]{正合结构}[exact structure]，\(\mathcal E\) 连同它称为\term[zhenghefanchou]{正合范畴}[exact category]。所选列称为可容许短正合列；其中 \(i,p\) 分别称为\term[kerongxudantai]{可容许单态射}[admissible monomorphism]与\term[kerongxumantai]{可容许满态射}[admissible epimorphism]。
\end{definition}

列两端可以添上 \(0\)，但不能据此认为范畴中每个态射都存在核和余核。这里的核、余核条件只针对选定的态射对。上述公理将分裂列直接列入；这是通常 Quillen 正合结构的一种等价写法。基础定义及例子可对照 Weibel《The K-book》\cite[Chapter II, Definition 7.0 and Examples 7.1.1--7.1.2]{Weibel2013KBook}。

\begin{proposition}\label{b4:prop:B-split-exact}
每个加性范畴 \(\mathcal E\) 都可配备只含分裂列的正合结构。
\end{proposition}
\begin{proof}
标准包含与投影 \(A\rightarrow A\oplus C\rightarrow C\) 互为核与余核；同构搬运不改变泛性质。两个标准分裂包含相继复合，仍同构于 \(A\rightarrow A\oplus C\oplus D\)；分裂投影的复合亦然。

对标准包含 \(A\rightarrow A\oplus C\) 及态射 \(f:A\rightarrow A'\)，其推出为 \(A'\oplus C\)，横向映射是 \(f\oplus1_C\)，新的包含仍为标准分裂包含。若两个从 \(A\oplus C,A'\) 出发的态射在 \(A\) 上相同，它们在 \(A'\) 与 \(C\) 上的分量唯一决定所需的推出态射，故确有泛性质。对偶地，标准投影 \(A\oplus C\rightarrow C\) 沿 \(g:C'\rightarrow C\) 的拉回为 \(A\oplus C'\)。全部公理成立。
\end{proof}

例如 \(\mathbb Z\)-有限自由模组成加性范畴，但乘 \(2:\mathbb Z\rightarrow\mathbb Z\) 在交换群范畴中的余核不是有限自由模。分裂正合结构不把它指定为可容许单态射。正合结构描述的是可以使用哪些短正合列，而非把一般单射一律重新命名。

\subsection{扩张封闭子范畴与可容许正合列}
\label{b4:subsec:B01:02}

\begin{proposition}\label{b4:prop:B-extension-closed}
设 \(\mathcal A\) 为交换范畴，\(\mathcal E\subset\mathcal A\) 为包含零对象、对同构封闭的全加性子范畴。若每个在 \(\mathcal A\) 中的短正合列
\(0\rightarrow A\rightarrow B\rightarrow C\rightarrow0\)
只要 \(A,C\in\mathcal E\) 就有 \(B\in\mathcal E\)，则把三项均在 \(\mathcal E\) 中的这些短正合列全部选取，便给出 \(\mathcal E\) 的正合结构。
\end{proposition}
\begin{proof}
核与余核的泛性质在全子范畴内仍成立；分裂列的中项为直和，属于 \(\mathcal E\)。设 \(A\rightarrow B\rightarrow D\) 为两个可容许单态射。交换范畴中的商对象满足
\[
0\longrightarrow B/A\longrightarrow D/A\longrightarrow D/B\longrightarrow0.
\]
两端属于 \(\mathcal E\)，由扩张封闭性，中间的 \(D/A\) 也属于 \(\mathcal E\)，所以复合可容许。对两个可容许满态射，复合的核是两个已知核的扩张，对偶地得到同一结论。

把 \(0\rightarrow A\rightarrow B\rightarrow C\rightarrow0\) 沿 \(A\rightarrow A'\) 推出，在 \(\mathcal A\) 中得到
\(0\rightarrow A'\rightarrow B'\rightarrow C\rightarrow0\)。
两端在 \(\mathcal E\)，所以 \(B'\) 在 \(\mathcal E\)。全性使这个推出仍满足 \(\mathcal E\) 内的泛性质。满态射的拉回同理，其核不变而商改为拉回的目标对象，因此中项也是 \(\mathcal E\) 中的扩张。
\end{proof}

有限生成投射左 \(R\)-模组成的范畴 \(R\text{-}\mathrm{proj}\) 是一个例子：投射模作商的短正合列分裂，中项为两端的有限直和。另一个例子是 \(\mathbb Z\)-有限生成无挠模，它也对扩张封闭。一般地，在交换范畴中保留一个扩张封闭对象族，往往比要求它对全部核、余核封闭更自然。

\ProseParagraphBreak
同一加性范畴可以有不同正合结构。在有限生成交换群上，通常的短正合结构包含
\[
0\longrightarrow\mathbb Z/p
\longrightarrow\mathbb Z/p^2
\longrightarrow\mathbb Z/p\longrightarrow0,
\]
而分裂结构不包含此列：中项有阶为 \(p^2\) 的元素，无法同构于两份 \(\mathbb Z/p\) 的直和。因此，写正合范畴时仅给对象和态射还不充分。

\begin{definition}\label{b4:def:B-exact-functor}
设 \(\mathcal E,\mathcal F\) 为指定了正合结构的加性范畴。若加性函子 \(F:\mathcal E\rightarrow\mathcal F\) 把每个可容许短正合列送到可容许短正合列，则称 \(F\) 为正合函子。
\end{definition}
任何加性函子都保持分裂列，所以对分裂正合结构自动正合；对较大的正合结构则须另行检查。例如 \(\operatorname{Hom}_{\mathbb Z}(\mathbb Z/p,-)\) 保持分裂列，却不把刚才的非分裂列变成满射结尾的短正合列。

\subsection{零调复形与有界导出范畴}
\label{b4:subsec:B01:03}

在一般正合范畴中，不能先假设全部上同调对象存在。因此应以指定短正合列描述“零调”。
\begin{definition}\label{b4:def:B-exact-acyclic}
设 \(\mathcal E\) 为正合范畴，\(X^\bullet\) 为其中的上链复形。若存在对象 \(Z^n\in\mathcal E\) 及可容许短正合列
\[
0\longrightarrow Z^n\xrightarrow{i^n}X^n
\xrightarrow{p^n}Z^{n+1}\longrightarrow0,
\qquad \mathrm{d}_X^n=i^{n+1}p^n,
\]
则称 \(X^\bullet\) 为关于该正合结构的零调复形。
\end{definition}

在交换范畴的通常正合结构中，这就是上同调全部为零。对扩张封闭全子范畴，还要求相应的 \(Z^n\) 仍属于所选子范畴。正合结构因而影响哪些复形将在导出范畴中变成零。

\begin{proposition}\label{b4:prop:B-split-acyclic-contractible}
设 \(\mathcal E\) 为幂等完备的加性范畴，配备分裂正合结构，\(X^\bullet\) 为 \(\mathcal E\) 中的上链复形。则 \(X^\bullet\) 零调，当且仅当它可缩，即存在 \(h^n:X^n\rightarrow X^{n-1}\) 满足
\(\mathrm{d}^{n-1}h^n+h^{n+1}\mathrm{d}^n=1_{X^n}\)。
\end{proposition}
\begin{proof}
若零调，每项可写为 \(X^n=Z^n\oplus Z^{n+1}\)，微分为 \((z,w)\mapsto(w,0)\)。令 \(h^n(z,w)=(0,z)\)，便得到所述等式。

反向设存在收缩 \(h\)。由 \(\mathrm{d}h\mathrm{d}=\mathrm{d}\)，每项上的 \(\mathrm{d}^{n-1}h^n\) 是幂等元，且与 \(h^{n+1}\mathrm{d}^n=1-\mathrm{d}^{n-1}h^n\) 互补。幂等完备性保证它们的像对象存在。记前者的像为 \(Z^n\)。微分把后一像同构地送到 \(Z^{n+1}\)，逆由限制后的 \(h^{n+1}\) 给出：两次复合在相应像上均由收缩等式为恒等。故 \(X^n\simeq Z^n\oplus Z^{n+1}\)，微分具有第一段的形式，得到分裂零调列。
\end{proof}

以下谈导出范畴时，额外假定 \(\mathcal E\) 幂等完备，以避免同伦同构下的零调封闭性问题。正合范畴的一般理论说明，零调有界复形组成 \(K^b(\mathcal E)\) 的厚三角子范畴 \(\operatorname{Ac}^b(\mathcal E)\)，于是可以定义
\[
D^b(\mathcal E)
=K^b(\mathcal E)/\operatorname{Ac}^b(\mathcal E).
\]
这里的商是 Verdier 商，等价于反演锥零调的复形映射。下面说明零调封闭性及商的构造；相关细节见 Bühler《Exact Categories》\cite[Definition 10.1, Proposition 10.14; Section 10.4]{Buehler2010ExactCategories}。
\begin{proofsketch}
零调复形的移位仍零调。设 \(X,Y\) 零调，\(f:X\to Y\) 为复形映射，以已选短正合列中的 \(Z^nX,Z^nY\) 表示相应余圈对象。锥的余圈对象可以由拉回构造：沿
\(-f:Z^{n+1}X\to Z^{n+1}Y\) 拉回可容许满态射 \(Y^n\to Z^{n+1}Y\)，得到
\[
 0\to Z^nY\to Z^n\operatorname{Cone}(f)
 \to Z^{n+1}X\to0.
\]
正合公理及其三乘三图追踪把这些列拼成锥的可容许短正合列。因此锥仍零调。直和因子的处理利用幂等完备性分裂这些余圈对象上的幂等元；同伦等价可在添加可缩复形后作同样处理。概要省略从同伦分裂改为逐项分裂的矩阵整理；这正是这里要求幂等完备而不任意删去该条件的原因。

令 \(S\) 为锥属于 \(\operatorname{Ac}^b(\mathcal E)\) 的态射。三角公理及刚证的封闭性使 \(S\) 含恒等、对复合封闭。把两个有共同定义域的映射放进同一个锥，可以形成同伦推出；若其中一边在 \(S\) 中，推出后的相应边之锥同构于原锥，仍在 \(S\) 中。这给出屋顶的共同分母。对差映射作同一锥构造，得到屋顶相等的消去条件。于是以 \(X\xleftarrow{s}U\to Y\)、\(s\in S\) 表示态射，再按共同细化识别，就得到三角局部化。

一个三角函子把 \(\operatorname{Ac}^b(\mathcal E)\) 送到零，当且仅当它把 \(S\) 中的态射送到同构。因此屋顶的像被唯一确定，得到 Verdier 商的泛性质。此构造解释了所写的 \(D^b(\mathcal E)\)，不需要在原正合范畴中先定义全部上同调对象。
\end{proofsketch}

在分裂情形，\cref{b4:prop:B-split-acyclic-contractible}说明被杀掉的复形在同伦范畴中已经为零，所以 \(D^b(\mathcal E)\simeq K^b(\mathcal E)\)。特别地，对 \(R\text{-}\mathrm{proj}\) 得到有界投射复形的同伦范畴。\cref{b4:prop:perfect-homotopy-model}进一步把它识别为 \(\operatorname{Perf}(R)\)：有界投射模型之间的导出态射就是同伦类，且每个完美对象按定义都有这样的模型。

% !TeX root = ../../Book4.tex
% 纲要标题：### B.2 \(K_0\) 与 Euler–Poincaré 可加性
% 纲要位置：工作记录/计划纲要.md，第 3577 行。
% 纲要细目（写作范围）：
% #### B.2.1 Grothendieck 群的定义与泛性质
% #### B.2.2 复形的 Euler 类与三角可加性
% #### B.2.3 有限生成投射模的类与算例

\section{\texorpdfstring{\(K_0\) 与 Euler–Poincaré 可加性}{K0 与 Euler–Poincaré 可加性}}
\label{b4:sec:B02}

维数把向量空间的短正合列变成整数加法。模的长度、秩以及复形的 Euler 特征也遵守类似规则。与其逐个选择数值，不如先构造一个交换群，使每个满足同一可加规律的不变量都从它得到。

\subsection{Grothendieck 群的定义与泛性质}
\label{b4:subsec:B02:01}

\begin{definition}\label{b4:def:B-grothendieck-group}
设 \(\mathcal E\) 为本质小的正合范畴。以对象同构类 \([X]\) 为生成元，以每个可容许短正合列
\(0\rightarrow A\rightarrow B\rightarrow C\rightarrow0\)
给出的
\[
[B]=[A]+[C]
\]
为关系构成的交换群，称为 \(\mathcal E\) 的 Grothendieck 群，记为 \(K_0(\mathcal E)\)。本质小性保证对象同构类组成集合。这把\cref{b4:def:grothendieck-group}中交换范畴的定义推广到指定的正合结构。
\end{definition}
\symidx{\(K_0(\mathcal E)\)：正合范畴的 Grothendieck 群}

取零列即得 \([0]=0\)，取分裂列即得 \([A\oplus C]=[A]+[C]\)。但若正合结构含非分裂列，还会加入新的关系。\([B]=[A]+[C]\) 一般并不意味着 \(B\simeq A\oplus C\)。

\begin{proposition}[可加不变量的泛性质]\label{b4:prop:B-k0-universal}
设 \(\mathcal E\) 为本质小正合范畴，\(M\) 为交换群，\(v\) 为从 \(\mathcal E\) 的对象同构类到 \(M\) 的函数。存在唯一群同态 \(\overline v:K_0(\mathcal E)\rightarrow M\)，使 \(\overline v\bigl([X]\bigr)=v(X)\) 对每个 \(X\in\mathcal E\) 成立，当且仅当 \(v(B)=v(A)+v(C)\) 对全部可容许短正合列 \(0\to A\to B\to C\to0\) 成立。
\end{proposition}
\begin{proof}
自由交换群的泛性质给出唯一的同态 \(\sum_i n_i[X_i]\mapsto\sum_i n_i\cdot v(X_i)\)。它下降到所定义的商群，当且仅当每个关系元 \([B]-[A]-[C]\) 被送到零，恰是所述条件。
\end{proof}

因而正合函子 \(F:\mathcal E\rightarrow\mathcal F\) 诱导群同态 \(K_0(F)\)，由 \([X]\mapsto\bigl[F(X)\bigr]\) 确定。对于分裂结构，\(K_0\) 正是对象同构类在直和下的交换幺半群的群完备化；若原范畴允许非分裂列，则须再除以这些列的关系。这里的定义和泛性质见 Weibel\cite[Chapter II, Definition 7.1]{Weibel2013KBook}。

\subsection{复形的 Euler 类与三角可加性}
\label{b4:subsec:B02:02}

\begin{proposition}[Euler–Poincaré 等式的重述]\label{b4:prop:B-euler-cohomology}
设 \(\mathcal A\) 为本质小交换范畴，\(X^\bullet\) 为其中的有界上链复形。则
\[
\sum_n(-1)^n[X^n]=\sum_n(-1)^n\bigl[H^n(X)\bigr]
\quad\text{在 }K_0(\mathcal A)\text{ 中成立}.
\]
\end{proposition}
\begin{proof}
这是\cref{b4:thm:euler-poincare}已经证明的结论。为说明下面推广时须检查什么，重述其两条短正合列。记 \(Z^n=\ker \mathrm{d}^n\)、\(B^n=\operatorname{im}\mathrm{d}^{n-1}\)，则
\[
0\rightarrow Z^n\rightarrow X^n\rightarrow B^{n+1}\rightarrow0,
\qquad
0\rightarrow B^n\rightarrow Z^n\rightarrow H^n(X)\rightarrow0
\]
给出 \([X^n]=[B^n]+\bigl[H^n(X)\bigr]+[B^{n+1}]\)。乘以 \((-1)^n\) 并求和，相邻的边界项抵消。由于复形有界，全部求和有限，没有收敛或换序问题。
\end{proof}

本命题在正合子范畴中使用时，还须保证上述边界、余圈和上同调对象都在该子范畴中，且这两列可容许。不能仅因每个 \(X^n\) 属于子范畴，就自动在它的 \(K_0\) 中写上同调的类。Weibel\cite[Chapter II, Proposition 7.5 and Corollary 7.5.1]{Weibel2013KBook}给出了保证这些条件的一类假设。

\begin{definition}\label{b4:def:B-triangulated-k0}
设 \(\mathcal T\) 为本质小三角范畴。以其对象同构类为生成元，并对每个好三角 \(A\rightarrow B\rightarrow C\rightarrow A[1]\) 加入关系 \([B]=[A]+[C]\)，所得交换群记为 \(K_0(\mathcal T)\)。
\end{definition}

对三角 \(X\rightarrow0\rightarrow X[1]\rightarrow X[1]\)，有 \(\bigl[X[1]\bigr]=-[X]\)。这解释了交错号与移位之间的关系，而不只是数值公式上的巧合。

\begin{theorem}[完美复形的 Euler 同构]\label{b4:thm:B-euler-isomorphism}
设 \(R\) 为含幺环，\(R\text{-}\mathrm{proj}\) 为有限生成投射左 \(R\)-模范畴，取通常的、亦即分裂的正合结构。令 \(K_0(R)=K_0(R\text{-}\mathrm{proj})\)。零次嵌入给出同构
\[
K_0(R)\ \simeq\ K_0\bigl(K^b(R\text{-}\mathrm{proj})\bigr)
\ \simeq\ K_0\bigl(\operatorname{Perf}(R)\bigr),
\]
第一个同构的逆由有界投射复形的 Euler 类
\(\chi(P^\bullet)=\sum_n(-1)^n[P^n]\)
给出。
\end{theorem}
\begin{proof}
先说明 Euler 类在同伦同构下不变。一个有界零调投射复形从最高非零次向下逐项分裂：最高项是投射模，前一微分到它的满射分裂，其核又为投射；重复便得到
\(P^n\simeq Z^n\oplus Z^{n+1}\)，其中各 \(Z^n\) 有限生成投射。因此该复形的 Euler 类为零。若 \(P\rightarrow Q\) 为同伦等价，其锥零调且仍有界投射，所以
\(\chi(Q)-\chi(P)=\chi\bigl(\operatorname{Cone}(f)\bigr)=0\)。

锥的 \(n\) 次项为 \(Q^n\oplus P^{n+1}\)，因此对锥三角有
\(\chi(Q)=\chi(P)+\chi\bigl(\operatorname{Cone}(f)\bigr)\)。
每个好三角同构于锥三角，所以 Euler 类诱导从三角 \(K_0\) 到 \(K_0(R)\) 的群同态。另一方面，投射模的可容许短正合列分裂；零次嵌入把它送到分裂三角，故也诱导反方向同态。Euler 类作用于零次投射模就是原来的类。

最后，对有界复形逐次取截去最低项后的子复形 \(F^{n+1}P\subset F^nP\)，逐次分裂短正合列给出好三角
\[
F^{n+1}P\longrightarrow F^nP\longrightarrow P^n[-n]
\longrightarrow F^{n+1}P[1].
\]
有限次使用三角关系及 \(\bigl[X[1]\bigr]=-[X]\)，得到
\([P^\bullet]=\sum_n(-1)^n[P^n]\)。
故另一方向复合亦为恒等。最后的同构来自\cref{b4:prop:perfect-homotopy-model}的三角等价：该等价全忠实且每个完美对象有有界投射模型，因而准确保持所用的好三角关系。
\end{proof}

例如对整数 \(m\ge2\)，\(\mathbb Z/m\) 作为完美 \(\mathbb Z\)-复形有两项模型 \(\mathbb Z\xrightarrow{m}\mathbb Z\)，故其 \(K_0\) 类为零，但这个导出对象并不为零。Grothendieck 群记录可加信息，不能完整分类对象。

\subsection{有限生成投射模的类与算例}
\label{b4:subsec:B02:03}

\begin{proposition}\label{b4:prop:B-basic-k0}
对任意域 \(k\)，有 \(K_0(k)\simeq\mathbb Z\)；并且 \(K_0(\mathbb Z)\simeq\mathbb Z\)。若 \(R=k_1\times\cdots\times k_r\) 为有限多个域的直积，则 \(K_0(R)\simeq\mathbb Z^r\)。
\end{proposition}
\begin{proof}
域上有限生成投射模就是有限维向量空间，同构类由维数分类，直和对应自然数相加，群完备化得到 \(\mathbb Z\)。整数上的有限生成投射模是有限自由模的直和项，因而是有限生成无挠交换群，由\cref{b4:thm:pid-free-submodule}关于主理想整环上自由模子模的结论，仍为自由模；有限生成性又使其秩有限；秩同样给出所需同构。

对域直积，令 \(e_i\) 为相应的中心幂等元。任一左 \(R\)-模分解为 \(\bigoplus_i e_iM\)，每个 \(e_iM\) 为 \(k_i\)-向量空间。有限生成投射模恰对应有限维向量空间的 \(r\)-元组：每个 \(Re_i\) 是 \(R\) 的投射直和项，而每个这样的元组都是有限多个 \(Re_i\) 的直和。维数向量因此给出 \(\mathbb N^r\)，群完备化为 \(\mathbb Z^r\)。
\end{proof}

若 \(R\) 为左 Noether 环，还可在全部有限生成左模的交换范畴上构造 \(G_0(R)\)。包含投射模的正合函子给出 \(K_0(R)\rightarrow G_0(R)\)，但二者使用的对象不同；不应只因记号相似便认为总同构。存在有限投射消解时，模的类可以用消解的 Euler 类表达；高投射维数的情形则需要进一步判断。

\ProseParagraphBreak
对 \(k\times k\)，\(K_0\) 已能区分两个投射直和项 \(k\times0\) 与 \(0\times k\)：它们的类分别为 \((1,0)\)、\((0,1)\)。只取底层 \(k\)-维数会把两者都送到 \(1\)，丢失了由两个中心幂等元确定的分量信息。

% !TeX root = ../../Book4.tex
% 纲要标题：### B.3 高阶代数 K 理论
% 纲要位置：工作记录/计划纲要.md，第 3882 行。
% 纲要细目（写作范围）：
% #### B.3.1 稳定一般线性群与一阶 K 群
% #### B.3.2 Q 构造与高阶 K 群
% #### B.3.3 局部化与高阶信息

\section{高阶代数 K 理论}
\label{b4:sec:B03}

\(K_0\) 将对象按可加关系组合，却没有直接保留对象的自同构。矩阵可逆变换提供了最初的高阶信息。再向上，关系之间的相容性需要用空间及其同伦群组织；这一节先把 \(K_1\) 算清，再说明一般 \(K_n\) 所依赖的准确构造。

\subsection{稳定一般线性群与一阶 K 群}
\label{b4:subsec:B03:01}

\begin{definition}\label{b4:def:B-k1}
设 \(R\) 为含幺环。通过 \(A\mapsto\operatorname{diag}(A,1)\) 把 \(\operatorname{GL}_n(R)\) 包含进 \(\operatorname{GL}_{n+1}(R)\)，所得并记为 \(\operatorname{GL}(R)\)，称为稳定一般线性群。其交换化
\[
K_1(R)\coloneqq
\operatorname{GL}(R)/\bigl[\operatorname{GL}(R),\operatorname{GL}(R)\bigr]
\]
称为 \(R\) 的\term[yijiedaishuKqun]{一阶代数 \(K\) 群}[algebraic \(K_1\)-group]。
\end{definition}
\symidx{\(\operatorname{GL}(R)\)：稳定一般线性群}
\symidx{\(K_1(R)\)：环的一阶代数 K 群}

交换子约定为 \([g,h]=ghg^{-1}h^{-1}\)。稳定化允许在矩阵右下角添加恒等块，而且确实可能消去固定大小交换化中的非零类。例如 \(\operatorname{GL}_2(\mathbb F_2)\) 的交换化为 \(C_2\)，而 \(K_1(\mathbb F_2)=\{1\}\)；其中二阶初等矩阵的非平凡类，在添加第三个指标后成为交换子的类。具体计算见\cref{b4:prop:B-stabilization-example}。定义可对照 Weibel\cite[Chapter III, Definition 1.1]{Weibel2013KBook}。

\begin{lemma}[稳定初等矩阵与交换子]\label{b4:lem:B-whitehead}
设 \(R\) 为含幺环，\(E(R)\) 为全部初等矩阵
\(e_{ij}(a)=I+aE_{ij}\)、\(i\ne j\)、\(a\in R\)
在 \(\operatorname{GL}(R)\) 中生成的子群。则
\[
E(R)=\bigl[\operatorname{GL}(R),\operatorname{GL}(R)\bigr].
\]
\end{lemma}
\begin{proof}
稳定化后总可选第三个指标 \(k\)。矩阵单位乘法给出
\(\bigl[e_{ik}(a),e_{kj}(1)\bigr]=e_{ij}(a)\)，故每个初等矩阵属于交换子群。

反向包含的目标是把任意交换子稳定后写成初等矩阵的乘积。为此，先把 \(g\) 与 \(g^{-1}\) 放在两个对角块上，证明 \(D(g)=\operatorname{diag}(g,g^{-1})\) 属于 \(E(R)\)；再用这些矩阵拼出交换子。要得到这样一对对角块，可以将带系数的交换矩阵与普通交换矩阵相乘，因此先对任意 \(g\in\operatorname{GL}_n(R)\) 写出
\[
w(g)=
\begin{pmatrix}I&g\\0&I\end{pmatrix}
\begin{pmatrix}I&0\\-g^{-1}&I\end{pmatrix}
\begin{pmatrix}I&g\\0&I\end{pmatrix}
=
\begin{pmatrix}0&g\\-g^{-1}&0\end{pmatrix}.
\]
两个形如上三角或下三角的块幺矩阵都是通常初等矩阵的乘积：在一个非对角块中的各矩阵单位两两乘积为零，逐个加入其系数即可。因此
\(D(g)=w(g)w(-I)=\operatorname{diag}(g,g^{-1})\)
属于 \(E(R)\)。对同样大小的 \(g,h\)，有
\[
D(g)D(h)D\bigl((hg)^{-1}\bigr)
=\operatorname{diag}\bigl([g,h],I\bigr).
\]
右下块为 \(g^{-1}h^{-1}hg=I\)，左上块为 \(ghg^{-1}h^{-1}\)。故每个交换子稳定后属于 \(E(R)\)，得到反向包含。
\end{proof}

这个证明给出 \(K_1(R)=\operatorname{GL}(R)/E(R)\)，并同时证明 \(E(R)\) 正规。无需事先断言每个固定大小的 \(E_n(R)\) 都在 \(\operatorname{GL}_n(R)\) 中正规；后一个断言对一般环并不成立。

\begin{theorem}\label{b4:thm:B-k1-field-integers}
对任意域 \(k\)，行列式诱导同构 \(K_1(k)\simeq k^\times\)。对整数环，有 \(K_1(\mathbb Z)\simeq\{1,-1\}\)。
\end{theorem}
\begin{proof}
行列式在稳定化下不变，且初等矩阵行列式为 \(1\)，故给出满同态 \(K_1(R)\rightarrow R^\times\)，其中 \(R=k\) 或 \(\mathbb Z\)。只需证明行列式为 \(1\) 的矩阵属于 \(E(R)\)。

域上可用行加法和带符号的行交换进行消元。带符号的交换矩阵
\(\left(\begin{smallmatrix}0&1\\-1&0\end{smallmatrix}\right)=w(1)\)
已由前一引理写成初等乘积。消元先得到非零对角元的上三角矩阵，再用行加法清除非对角项，得到行列式为 \(1\) 的对角矩阵。每个块 \(\operatorname{diag}(u,u^{-1})=D(u)\) 为初等乘积，因此任意乘积为 \(1\) 的对角元也可相继消去。

整数情形中，行列式为 \(1\) 使第一列元素的最大公因数为 \(1\)。两行间的带余除法与带符号交换，就是 Euclid 算法\footnote{欧几里得（\OriginalName{Εὐκλείδης}，约前325-约前265），古希腊数学家，著有《几何原本》，系统整理了几何学与数论。}，可把这一列化为 \((1,0,\ldots,0)^t\)；若先得到 \(-1\)，用 \(D(-1)\) 同时改变两行符号。然后用列加法清掉第一行剩余元素。余下方块仍为行列式 \(1\) 的整数矩阵，按大小归纳得到恒等矩阵。列加法也是右乘初等矩阵，故原矩阵仍在 \(E(\mathbb Z)\) 中。大小为 \(1\) 的情形直接成立；需要两行时可先稳定化一次。两种情形的行列式核均为初等子群，结合\cref{b4:lem:B-whitehead}完成证明。
\end{proof}

域上的 \(K_1(k)=k^\times\) 保留了乘法群的信息，而 \(K_0(k)=\mathbb Z\) 只反映维数。一般交换环还可能有非零的行列式核
\(\operatorname{SK}_1(R)=\ker\bigl(K_1(R)\rightarrow R^\times\bigr)\)，所以不能把域上的公式无条件推广。非零核的例子见 Weibel\cite[Chapter III, Example 1.5.4]{Weibel2013KBook}。

\subsection{Q 构造与高阶 K 群}
\label{b4:subsec:B03:02}

设 \(\mathcal E\) 为小正合范畴。Quillen 的 \(Q\) 构造保留对象，却以子商数据作为态射：从 \(A\) 到 \(B\) 的一张态射，指定 \(B\) 的一个可容许子对象，以及该子对象的一个可容许商 \(A\)。这不是原范畴中的一张 \(A\to B\) 映射。
\begin{definition}\label{b4:def:B-q-construction}
设 \(\mathcal E\) 为小正合范畴。构造范畴 \(Q\mathcal E\)，其对象与 \(\mathcal E\) 相同，从 \(A\) 到 \(B\) 的态射为图式
\[
A\xleftarrow{p}M\xrightarrow{i}B
\]
的同构类，其中 \(p\) 为可容许满态射，\(i\) 为可容许单态射。中间对象的同构必须使两端映射均相容。与 \(B\xleftarrow{q}N\xrightarrow{j}C\) 的复合由拉回 \(M\times_BN\) 给出：
\[
A\longleftarrow M\times_BN\longrightarrow C.
\]
恒等态射由两条恒等映射 \(A\leftarrow A\to A\) 表示。这个构造称为 Quillen 的 \(Q\) 构造。
\end{definition}
\symidx{\(Q\mathcal E\)：正合范畴的 Quillen \(Q\) 构造}

在模或向量空间中，用 \(i\) 把 \(M\) 识别为 \(B\) 的子对象，则复合中的拉回就是 \(q^{-1}(M)\subseteq N\)。我们先在 \(C\) 的子对象 \(N\) 中找出映入 \(M\) 的部分，再经 \(q\) 和 \(p\) 将它映到 \(A\)。因此新图式的左端仍是一个商，右端仍是一个子对象。特别地，从 \(0\) 到向量空间 \(B\) 的 \(Q\)-态射恰好对应 \(B\) 的子空间；当 \(B=\mathbb F_q^2\) 时共有 \(q+3\) 张，而非唯一一张。这个对应及计数见\cref{b4:prop:B-q-zero-subspaces}。

\ProseParagraphBreak
正合公理保证这个拉回存在，投影到 \(M\) 为可容许满态射。另取 \(i:M\rightarrow B\) 的可容许余核 \(B\rightarrow D\)，则复合 \(N\xrightarrow{q}B\rightarrow D\) 仍为可容许满态射。其核正是 \(M\times_BN\rightarrow N\)：一个到 \(N\) 的态射在复合到 \(D\) 后为零，当且仅当其到 \(B\) 的像经由 \(M\)，再由拉回泛性质得到唯一分解。因此该投影为可容许单态射，两端复合仍属所定类型。重复拉回的泛性质给出典范同构，通过中间对象同构类取商后，复合严格结合。恒等态射由 \(A\leftarrow A\rightarrow A\) 两个恒等映射表示。

\ProseParagraphBreak
范畴的神经 \(N\mathcal C\) 把可复合的 \(n\) 个态射作为 \(n\)-单形；删去端点或复合相邻态射给出面映射，插入恒等给出退化映射。其几何实现记为 \(B\mathcal C\)。这个空间同时反映对象、态射及复合的相容关系。本节假定读者熟悉同伦群等代数拓扑基础，不在附录中重述这些理论。
\symidx{\(N\mathcal C,\ B\mathcal C\)：范畴的神经与分类空间}

\ProseParagraphBreak
在 \(Q\mathcal E\) 中，每个对象 \(A\) 都有从 \(0\) 来的态射，所以 \(BQ\mathcal E\) 连通；\(\pi_0\) 无法区分对象的可加类。不过，从 \(0\) 到 \(A\) 有两张特别的态射
\[
 \alpha_A:\quad 0\longleftarrow0\longrightarrow A,
 \qquad
 \beta_A:\quad 0\longleftarrow A\xrightarrow{1_A}A.
\]
它们在分类空间中给出两条从 \(0\) 到 \(A\) 的路径。以下采用先沿 \(\beta_A\) 前进、再沿 \(\alpha_A\) 反向返回的环路方向，对应 \(A\) 的正可加类；反过来走给出它的负类。对象的可加信息因而可以保留在基本群中，这正是下面将 \(K_0\) 放在 \(\pi_1\) 的原因。

\begin{definition}\label{b4:def:B-higher-k}
设 \(\mathcal E\) 为小正合范畴，以零对象作为 \(BQ\mathcal E\) 的基点。对 \(n\ge0\)，定义
\[
K_n(\mathcal E)=\pi_{n+1}(BQ\mathcal E,0).
\]
对本质小范畴先选取小骨架；等价范畴的神经同伦等价保证此选择不影响所得群。对含幺环 \(R\)，以 \(R\text{-}\mathrm{proj}\) 表示有限生成投射左 \(R\)-模范畴，取分裂正合结构，并定义 \(K_n(R)=K_n(R\text{-}\mathrm{proj})\)。
\end{definition}
\symidx{\(K_n(\mathcal E),K_n(R)\)：正合范畴与环的代数 \(K\) 群}

这些定义与前文的低阶定义相容：
\(\pi_1(BQ\mathcal E)\simeq K_0(\mathcal E)\)，且环上的 \(K_1\) 同构于稳定一般线性群的交换化。由此，前一小节的矩阵计算也给出了本节定义的 \(K_1\)。比较结果见 Weibel\cite[Chapter IV, Proposition 6.2; Section 7]{Weibel2013KBook}。
前一种比较直接计算 \(BQ\mathcal E\) 的基本群；后一种比较使用其环路空间的基点连通分支，将 \(\pi_2(BQ\mathcal E)\) 转成这一分支的基本群，再以加号构造计算它。群完成和加号构造的空间比较仍需所引的拓扑论证。
\begin{proofsketch}
范畴神经的一个单形就是一串可复合态射，因此其基本群可以用边和二单形的复合关系计算。取上述 \(\alpha_A\) 为连接各顶点的树，\(\beta_A\) 后接 \(\alpha_A\) 反向路径的环路类记为 \([A]\)。每个可容许短正合列 \(0\to A\to B\to C\to0\) 给出关系 \([B]=[A]+[C]\)。反过来，每条边都表示一个可容许子商；若由 \(A\leftarrow M\to B\) 表示，其相对于上述树的环路类由 \(\ker(M\to A)\) 的类给出。特别地，\(\alpha_B\) 是树边，给出零类；\(\beta_B\) 给出 \([B]\)。拉回计算说明，复合两条边时相应核的类相加。二单形关系进一步给出可容许短正合列的加法关系；从全部子商关系归约到这些关系的图式追踪见所引低阶比较的证明。直和的两个次序同构，使生成元彼此交换，得到第一项。

对 \(\mathcal E=R\text{-}\mathrm{proj}\)，全部正合列分裂。令 \(\mathcal S\) 为投射模及同构组成的范畴，以直和为运算；再用短正合列组成辅助范畴，按末项投影到 \(Q\mathcal E\)。在两边同时添加一个直和因子，并把此操作形式地变为可逆，可以比较它的纤维与 \(\mathcal S\) 的群完成。拉回短正合列给出纤维之间的比较；分裂性使这些比较为同伦等价。因此群完成空间与 \(\Omega BQ\mathcal E\) 同伦等价。这一步的范畴定义与纤维比较见\cite[Chapter IV, Definition 7.3; Example 7.4; Lemma 7.5; Proposition 7.6; Theorem 7.8]{Weibel2013KBook}，这里省略神经实现与群完成的技术验证。

每个有限生成投射左模都有补模，添加补模后成为有限自由左模。把自由左模 \({}_RR^n\) 的元素写成行向量，则每个自同构唯一写成 \(r_A:x\mapsto xA\)，其中 \(A\in\operatorname{GL}_n(R)\)。由于 \(r_A\circ r_B=r_{BA}\)，其自同构群是 \(\operatorname{GL}_n(R)^{\mathrm{op}}\)。逆矩阵映射 \(A\mapsto A^{-1}\) 给出从这个反群到 \(\operatorname{GL}_n(R)\) 的群同构，与块稳定化及块直和相容，并把初等矩阵 \(e_{ij}(a)\) 送到 \(e_{ij}(-a)\)。因此群完成与加号构造的比较给出
\[
\begin{aligned}
 \pi_0\bigl(\Omega BQ\mathcal E\bigr)&\cong K_0(R),\\
 \bigl(\Omega BQ\mathcal E\bigr)_0
 &\simeq B\bigl(\operatorname{GL}(R)^{\mathrm{op}}\bigr)^+
 \simeq B\operatorname{GL}(R)^+.
\end{aligned}
\]
其中下标 \(0\) 表示基点所在的连通分支。加号构造杀掉初等子群 \(E(R)\)，并保持同调；\cref{b4:lem:B-whitehead}给出 \(E(R)=[\operatorname{GL}(R),\operatorname{GL}(R)]\)，其证明中的初等交换子公式还表明 \(E(R)\) 自身是完美群。

逆矩阵识别把自同构 \(r_A\) 的环路类送到 \([A^{-1}]\)。群完成的零分支是由直和运算给出的群状 \(H\)-空间，它的同伦逆在 \(\pi_1\) 上再次取逆。将上述比较与这项同伦逆复合，便得到把 \([r_A]\) 送到 \([A]\) 的比较；以下使用这一约定。于是
\[
 \pi_2(BQ\mathcal E)\cong\pi_1\bigl((\Omega BQ\mathcal E)_0\bigr)
 \cong\operatorname{GL}(R)/E(R)=K_1(R).
\]
群完成与加号构造的空间比较见\cite[Chapter IV, Theorems 4.10 and 7.1; Corollary 7.2; Exercise 7.9]{Weibel2013KBook}。该书的环版本采用右模；这里先把对称幺半范畴的群完成及分裂正合范畴的一般比较用于左模，再通过反群的逆矩阵识别得到同一加号空间。
\end{proofsketch}

\subsection{局部化与高阶信息}
\label{b4:subsec:B03:03}

高阶 \(K\) 群使零阶的可加关系延伸为长正合列。以下给出一个有明确适用范围的代表性结果。

\ProseParagraphBreak
先说明式中商范畴的意义。交换范畴 \(\mathcal A\) 的 Serre 子范畴 \(\mathcal B\) 对子对象、商对象与扩张封闭。Serre 商 \(\mathcal A/\mathcal B\) 保留原来的对象，并把核、余核都属于 \(\mathcal B\) 的态射变成同构；其中 \(A\) 成为零对象，当且仅当 \(A\in\mathcal B\)。于是包含函子记录被消去的部分，商函子记录消去后仍可见的部分，长列中的三组 \(K\) 群分别属于这三个正合范畴。
\symidx{\(\mathcal A/\mathcal B\)：交换范畴关于 Serre 子范畴的商}

\ProseParagraphBreak
例如在有限生成交换群范畴中消去有限群，有限指数的包含 \(n\mathbb Z\hookrightarrow\mathbb Z\) 的核为零、余核有限，因而在商范畴中可逆。屋顶
\[
 \mathbb Z\longleftarrow n\mathbb Z\longrightarrow\mathbb Z,
 \qquad nm\longmapsto m\quad(n\ge1),
\]
便表示乘以 \(1/n\)。这解释了为何该商范畴的态射会出现有理系数，而不是仅把每个对象的有限挠子群删去后仍保留原有态射。

\begin{theorem}[Quillen 局部化]\label{b4:thm:B-quillen-localization}
设 \(\mathcal A\) 为小交换范畴，\(\mathcal B\subset\mathcal A\) 为 Serre 子范畴，即对子对象、商对象与扩张封闭的全子范畴。则包含与 Serre 商函子给出长正合列，其中中间段的 \(n\ge1\)：
\[
\begin{aligned}
\cdots&\longrightarrow K_{n+1}(\mathcal A/\mathcal B)
 \xrightarrow{\partial}K_n(\mathcal B)
 \longrightarrow K_n(\mathcal A)\\
&\longrightarrow K_n(\mathcal A/\mathcal B)
 \xrightarrow{\partial}K_{n-1}(\mathcal B)\longrightarrow\cdots\\
&\longrightarrow K_0(\mathcal B)\longrightarrow K_0(\mathcal A)
 \longrightarrow K_0(\mathcal A/\mathcal B)\longrightarrow0.
\end{aligned}
\]
本质小情形通过小骨架作同样解释。
\end{theorem}
\begin{proofsketch}
写 \(q:\mathcal A\to\mathcal A/\mathcal B\)，考虑
\(Qq:Q\mathcal A\to Q(\mathcal A/\mathcal B)\)。对商范畴的对象 \(L\)，令 \(\mathcal D_L\) 的对象为 \(A\in\mathcal A\) 连同 \(Q(\mathcal A/\mathcal B)\) 中的态射 \(L\to q(A)\)。这是计算该函子同伦纤维的逗号范畴。其全子范畴 \(\mathcal F_L\) 只保留所给态射为同构的对象。

商范畴中的子商可以由原范畴中的层 \(A_1\subseteq A_2\subseteq A\) 表示。同一个子商的两个表示可通过对子对象取交与和得到共同细化；比较这些表示的映射，其核与余核属于 \(\mathcal B\)，Serre 封闭性使细化中相应的差商仍属于 \(\mathcal B\)。固定一个子商表示后，所有细化所成的范畴是滤过的，因而其神经可缩。Quillen 的定理 A 遂给出 \(B\mathcal F_L\simeq B\mathcal D_L\)。当 \(L=0\) 时，\(q(A)\cong0\) 恰好等价于 \(A\in\mathcal B\)，故
\[
 B\mathcal D_0\simeq B\mathcal F_0\simeq BQ\mathcal B.
\]

还须证明商范畴中的态射引起 \(\mathcal D_L\) 之间的同伦等价。为此，对 \(N\in\mathcal A\) 考察所有在商范畴中成为同构的映射 \(A\to N\)。先取像，便把它化成满射；核属于 \(\mathcal B\)。将核作可容许扩张并用拉回、推出改变表示，可以比较这个辅助范畴与 \(Q\mathcal B\)。所有像 \(N'\subseteq N\)、\(N/N'\in\mathcal B\) 有共同细化，因而取不同的像不改变其同伦型。由此得到所需的换基同伦等价。

Quillen 的定理 B 现在识别出同伦纤维列
\[
 BQ\mathcal B\longrightarrow BQ\mathcal A
 \longrightarrow BQ(\mathcal A/\mathcal B).
\]
同伦群长正合列的典型片段为
\[
 \pi_{n+1}(BQ\mathcal A)
 \longrightarrow\pi_{n+1}\bigl(BQ(\mathcal A/\mathcal B)\bigr)
 \xrightarrow{\partial}\pi_n(BQ\mathcal B),
 \qquad n\ge1.
\]
代入 \(K_n=\pi_{n+1}BQ\)，右端便是 \(K_{n-1}(\mathcal B)\)：连接映射降低一次 \(K\) 次数，来自同伦长正合列的次数变化。由于 \(BQ\mathcal B\) 连通，其 \(\pi_0\) 为零，故末尾的 \(K_0(\mathcal A)\to K_0(\mathcal A/\mathcal B)\) 满射。概要省略定理 A、B 的拓扑证明和上述辅助范畴的全部自然变换；前者见\cite[Chapter IV, Theorems 3.7--3.8]{Weibel2013KBook}，后者按\cite[Chapter V, Theorem 5.1, Claims 5.1.2--5.1.7]{Weibel2013KBook}逐项验证。
\end{proofsketch}

令 \(\mathcal A\) 为有限生成交换群，\(\mathcal B\) 为有限交换群。张量 \(\mathbb Q\) 杀掉且仅杀掉有限群，Serre 商等价于有限维 \(\mathbb Q\)-向量空间范畴。态射上，这个识别来自清除分母：有理线性映射在某个有限指数子群上成为整数线性映射，而有限指数的变化在商范畴中成为同构。有限交换群的合成因子都是素阶循环群，所以短正合关系把每个 \([M]\) 写成 \([\mathbb Z/p\mathbb Z]\) 的有限整数和。另一方面，短正合列中的群阶相乘，故各 \(p\)-主分量的长度 \(v_p(|M|)\) 诱导群同态
\[
K_0(\mathcal B)\longrightarrow\bigoplus_{p\;\text{为素数}}\mathbb Z,
\qquad [M]\longmapsto\bigl(v_p(|M|)\bigr)_p.
\]
它把 \([\mathbb Z/p\mathbb Z]\) 送到第 \(p\) 个基向量，与将该基向量送回 \([\mathbb Z/p\mathbb Z]\) 的群同态互为逆，因而得到 \(K_0(\mathcal B)\simeq\bigoplus_p\mathbb Z\)，没有额外的生成元关系。在 \(K_0(\mathcal A)\) 中，短正合列
\(0\rightarrow\mathbb Z\xrightarrow{p}\mathbb Z\rightarrow\mathbb Z/p\mathbb Z\rightarrow0\)
使这些生成元全为零。有限生成交换群的结构定理因而给出
\(K_0(\mathcal A)\simeq\mathbb Z\)，由秩识别。

\ProseParagraphBreak
连接映射如何产生有限群的类，可以先看素数 \(p\)。整数映射 \(p:\mathbb Z\to\mathbb Z\) 的核为零、余核为 \(\mathbb Z/p\mathbb Z\)，所以它在 Serre 商中成为同构，且对应 \(\mathbb Q\) 上乘 \(p\) 的自同构。在原范畴中，这个提升未能成为同构的部分正是其余核。按取余核类减核类的符号约定，连接映射给出
\[
 \begin{aligned}
 \partial(p)&=[\mathbb Z/p\mathbb Z],\\
 \partial(p^m)&=m[\mathbb Z/p\mathbb Z],\\
 \partial(p^{-m})&=-m[\mathbb Z/p\mathbb Z]\quad(m\ge1).
 \end{aligned}
\]
后一组等式来自 \(\partial\) 从乘法群到加法群的同态性质。乘 \(-1\) 的整数映射本来就是自同构，核与余核都为零，故 \(\partial(-1)=0\)。该余核减核的公式见 Weibel\cite[Chapter V, Application 6.1 and Example 6.1.2]{Weibel2013KBook}。将任意非零有理数分解为符号与有限多个素数幂，就得到素数赋值
\[
\mathbb Q^\times\longrightarrow\bigoplus_p\mathbb Z,
\qquad a\longmapsto\bigl(v_p(a)\bigr)_p.
\]
仅从算术也能直接验证
\[
1\longrightarrow\{1,-1\}\longrightarrow\mathbb Q^\times
\longrightarrow\bigoplus_p\mathbb Z\longrightarrow0
\]
正合：素因子分解给出核与满射性。这说明在零阶可加关系中消失的有限群类，会作为高一阶连接映射的值重新出现。

\begin{proposition}\label{b4:prop:B-stabilization-example}
在域 \(\mathbb F_2\) 上，\(\operatorname{GL}_2(\mathbb F_2)\simeq S_3\)，它的交换化为 \(C_2\)，但稳定交换化 \(K_1(\mathbb F_2)\) 是平凡群。二阶初等矩阵 \(e_{12}(1)\) 在前一个交换化中给出非平凡类，添加一个坐标后成为交换子。
\end{proposition}
\begin{proof}
可逆矩阵置换 \(\mathbb F_2^2\) 的三个非零向量。如果某矩阵固定它们，就特别固定一组基，因而是恒等矩阵，所以这个作用忠实。可逆矩阵的第一列有 \(3\) 种选择，第二列有 \(2\) 种选择，总数为 \(6\)，故所得到 \(S_3\) 的单同态是同构。\(S_3\) 的符号同态核为 \(A_3\)，而两个换位的交换子给出一个三轮换，故其交换子群恰为 \(A_3\)，交换化为 \(C_2\)。矩阵 \(e_{12}(1)\) 固定 \((1,0)^t\)，交换 \((0,1)^t\) 与 \((1,1)^t\)，所以对应换位并给出非平凡交换化类。稳定化后有
\[
 e_{12}(1)=\bigl[e_{13}(1),e_{32}(1)\bigr],
\]
故这个类在稳定交换化中变为单位类。最后，\cref{b4:thm:B-k1-field-integers}给出 \(K_1(\mathbb F_2)\simeq\mathbb F_2^\times=\{1\}\)。
\end{proof}

\begin{proposition}\label{b4:prop:B-q-zero-subspaces}
设 \(q\) 为素数幂，\(\mathcal E\) 为以 \(\mathbb F_q^n\)（\(n\ge0\)）为对象的有限维向量空间范畴的标准小骨架，取通常的正合结构。对 \(V\in\mathcal E\)，\(Q\mathcal E\) 中从 \(0\) 到 \(V\) 的态射与 \(V\) 的子空间一一对应。特别地，
\[
 \#\operatorname{Hom}_{Q\mathcal E}(0,\mathbb F_q^2)=q+3.
\]
原范畴的零对象 \(0\) 在 \(Q\mathcal E\) 中既不是始对象，也不是终对象。
\end{proposition}
\begin{proof}
从 \(0\) 到 \(V\) 的态射由 \(0\leftarrow M\xrightarrow{i}V\) 表示。到零对象的映射自动满射，而单射 \(i\) 把 \(M\) 识别为 \(V\) 中的子空间 \(i(M)\)。两张图式等价当且仅当它们有相同的像：若像相同，经各自单射识别便得到使图式相容的中间对象同构；若图式等价，则像当然相同。每个子空间又能选取骨架中的同构代表，故得到所述对应。

二维空间的零子空间与全空间各一个；一维子空间的非零向量彼此不交，每个含 \(q-1\) 个非零向量，故一维子空间有 \((q^2-1)/(q-1)=q+1\) 个。总数为 \(q+3\)，所以 \(0\) 不是始对象。若 \(V\ne0\)，一张 \(V\to0\) 的 \(Q\)-态射需要单射 \(M\to0\)，迫使 \(M=0\)，却又要求 \(M\to V\) 满射，矛盾。因此 \(0\) 也不是终对象。
\end{proof}

% !TeX root = ../../Book4.tex
% 纲要标题：### B.4 模型范畴与稳定无穷范畴
% 纲要位置：工作记录/计划纲要.md，第 3593 行。
% 纲要细目（写作范围）：
% #### B.4.1 模型结构与导出函子
% #### B.4.2 稳定无穷范畴与同伦相干性
% ---

\section{模型范畴与稳定无穷范畴}
\label{b4:sec:B04}

导出范畴把拟同构变成同构；具体消解则为计算保留了可操作的代表。模型范畴把这两方面组织成统一公理。若还希望保留态射之间的同伦及其高阶相容性，就需要比普通同伦范畴更丰富的结构。这里通过复形说明两种语言各自增加了什么。

\subsection{模型结构与导出函子}
\label{b4:subsec:B04:01}

\begin{definition}\label{b4:def:B-model-category}
设 \(\mathcal M\) 为有全部小极限和小余极限的范畴，指定三类态射，分别称为弱等价、余纤维化和纤维化。若它们满足：
\begin{enumerate}
\item 弱等价满足二取三性质；三类态射均对同构和收缩封闭。
\item 余纤维化对同时为弱等价的纤维化有左提升性质；同时为弱等价的余纤维化对全部纤维化有左提升性质。
\item 每个态射可函子性地分解为“余纤维化后接平凡纤维化”，也可分解为“平凡余纤维化后接纤维化”。
\end{enumerate}
则称这一指定为\term[moxingjiegou]{模型结构}[model structure]。这里“平凡”表示同时为弱等价；左提升性质是指每个以这两条态射为左边和右边的交换方块都有对角填充。带此结构的范畴称为\term[moxingfanchou]{模型范畴}[model category]。
\end{definition}

这是采用函子性分解的通常版本，见 Weibel\cite[Chapter V, Definition 10.4]{Weibel2013KBook}。从始对象到 \(X\) 为余纤维化时，称 \(X\) 余纤维；从 \(X\) 到终对象为纤维化时，称 \(X\) 纤维。两种分解分别产生余纤维替换 \(QX\rightarrow X\) 与纤维替换 \(X\rightarrow RX\)。同伦范畴 \(\operatorname{Ho}(\mathcal M)\) 则通过反演弱等价得到。

\begin{example}\label{b4:exm:B-field-complex-model}
设 \(k\) 为域。在全部 \(k\)-向量空间上链复形组成的范畴 \(\operatorname{Ch}(k)\) 上，可取拟同构为弱等价、逐次单射为余纤维化、逐次满射为纤维化。这给出模型结构，且每个对象同时余纤维、纤维。
\end{example}
\begin{proof}
范畴的极限、余极限逐次存在；二取三由上同调给出，收缩封闭性由单射、满射及同构的收缩性质给出。下面核验其余两项。

域上的零调复形都可缩：逐次选择余圈的补空间，得到 \(X^n=Z^n\oplus Z^{n+1}\) 及微分 \((z,w)\mapsto(w,0)\)。它是两项恒等复形的直和，也是在每个次数仅有两项贡献的同一族的直积。两项恒等复形在 \(\operatorname{Ch}(k)\) 中既投射又内射：从它出发的复形映射由一个次数上的线性映射决定，映入它的复形映射也由一个次数上的线性映射决定，而向量空间范畴中的 Hom 正合。因此零调复形在复形范畴中也既投射又内射。

若满射 \(p:E\rightarrow F\) 为拟同构，其核 \(K\) 零调而内射，所以 \(E\simeq K\oplus F\)。对给定的提升问题，指向 \(F\) 的分量已经确定，指向 \(K\) 的分量可沿左边的逐次单射延拓，因为 \(K\) 在复形范畴内内射。这证明第一种提升。若左边的单射为拟同构，其余核零调而投射，短正合列分裂；在这个余核上沿右边满射提升，便得到另一种提升。

最后，记 \(C(X)=\operatorname{Cone}(1_X)\)，它可缩，并有逐次单射 \(j_X:X\rightarrow C(X)\)。任意复形映射 \(f:X\rightarrow Y\) 有分解
\[
X\xrightarrow{(f,j_X)}Y\oplus C(X)\xrightarrow{\mathrm{pr}_Y}Y,
\]
第一箭头逐次单射，第二箭头为逐次满射拟同构。又 \(C(Y)[-1]\) 有逐次满射 \(q_Y:C(Y)[-1]\rightarrow Y\)，在锥坐标上取第二分量。于是还有
\[
X\longrightarrow X\oplus C(Y)[-1]
\xrightarrow{(f,q_Y)}Y,
\]
第一箭头为标准包含，是逐次单射拟同构，第二箭头逐次满射。这些锥构造对交换方块函子性地作用，所以分解满足要求。最后，\(0\rightarrow X\) 总是逐次单射，\(X\rightarrow0\) 总是逐次满射。
\end{proof}

这个例子的简便性来自底域条件。一般环上的零调复形未必可缩，不能仍把所有逐次单射、满射原样作为同一套模型结构。投射消解与内射消解的选择正是在此重新发挥作用。

\ProseParagraphBreak
设 \(\mathcal M,\mathcal N\) 为模型范畴，\(L:\mathcal M\rightleftarrows\mathcal N:U\) 为伴随对。若左伴随保持余纤维化及平凡余纤维化，则称其为\term[Quillenbansui]{Quillen 伴随}[Quillen adjunction]。模型范畴的一般定理给出同伦范畴上的导出伴随，取余纤维替换 \(X^{\mathrm{cof}}\) 与纤维替换 \(Y^{\mathrm{fib}}\) 后写成
\[
\mathbf L L(X)=L(X^{\mathrm{cof}}),\qquad
\mathbf R U(Y)=U(Y^{\mathrm{fib}}).
\]
这概括了“先选择适当消解再施行函子”的做法。
\begin{proofsketch}
先说明 \(L\) 把余纤维对象之间的弱等价送到弱等价。对这样的 \(f:X\to Y\)，把 \((f,1_Y):X\amalg Y\to Y\) 分解为余纤维化 \(j:X\amalg Y\to Z\) 后接平凡纤维化 \(p:Z\to Y\)。两项包含后所得 \(j_X,j_Y\) 都是余纤维化；由 \(pj_X=f\)、\(pj_Y=1_Y\) 和二取三性质，它们也都是弱等价。故 \(L(j_X),L(j_Y)\) 为弱等价。又 \(L(p)L(j_Y)=1\)，二取三先给出 \(L(p)\) 为弱等价，再给出 \(L(f)\) 为弱等价。对偶论证说明 \(U\) 保持纤维对象之间的弱等价。

若 \(QX\to X\)、\(Q'X\to X\) 是两个余纤维替换，提升性质给出它们之间的比较，比较仍为弱等价；因此施行 \(L\) 后在同伦范畴中同构。选择函子性替换使对象公式同时定义态射；其对原弱等价的保持由上一步得到。右导出的论证对偶。

取余纤维 \(X\) 和纤维 \(Y\)。普通伴随给出 \(\mathcal N(LX,Y)\cong\mathcal M(X,UY)\)。把一个柱对象施行 \(L\)，由左伴随保持余积、余纤维化和刚证的弱等价，仍得到 \(LX\) 的柱对象。把一个路径对象施行 \(U\)，由右伴随保持积、纤维化及纤维对象间弱等价，仍得到 \(UY\) 的路径对象。于是此双射同时保持两个方向的同伦关系，下降为
\[
 \operatorname{Ho}(\mathcal N)(\mathbf LL(X),Y)
 \cong\operatorname{Ho}(\mathcal M)(X,\mathbf RU(Y)).
\]
这里使用余纤维对象到纤维对象的同伦类计算同伦范畴态射；其依据是把对象同时替换为余纤维、纤维对象，提升公理使此类对象间的弱等价有同伦逆，故局部化不会再增加新的态射。概要省略一般模型范畴中左同伦与右同伦一致性的逐项提升图。上面的域上复形模型中，这些比较就是已经计算的链同伦。
\end{proofsketch}

\subsection{稳定无穷范畴与同伦相干性}
\label{b4:subsec:B04:02}

本卷的 DG 范畴为每对复形指定一个总 Hom 复形。由\cref{b4:prop:dg-hom-complex}，其零次上同调给出普通同伦态射，其他次数则给出移位态射。只留下同伦态射并不会删掉三角范畴中连接映射的同伦类，却会丢掉具体同伦以及构造连接映射时的高阶相干数据。

\ProseParagraphBreak
无穷范畴的一种理解是同时保留态射、态射之间的可逆同伦，以及这些同伦之间的高阶相容性。精确模型可以采用满足内角填充条件的单纯集，称为拟范畴；这里不展开整套单纯模型，而在已给定无穷范畴的前提下说明稳定条件。相关语言及模型比较见 Blumberg、Gepner、Tabuada 的《A universal characterization of higher algebraic K-theory》\cite[Section 2.2, pp. 747--750]{BlumbergGepnerTabuada2013UniversalKTheory}。

\begin{definition}\label{b4:def:B-stable-infinity}
设 \(\mathcal C\) 为无穷范畴。若它具有零对象、全部有限极限与有限余极限，并且其中一个方块为推出方块当且仅当它为拉回方块，则称 \(\mathcal C\) 为稳定高阶范畴，也称稳定无穷范畴。\termidx[wendingwuqiongfanchou]{稳定无穷范畴}这里的极限与余极限均在无穷范畴的意义下理解，其泛性质比较的是映射空间的等价。
\end{definition}

这个条件使同伦纤维与同伦余纤维只差一次移位。稳定无穷范畴的同伦范畴具有三角结构，好三角来自余纤维序列，参见\cite[Definition 2.12 and the following paragraph, p. 749]{BlumbergGepnerTabuada2013UniversalKTheory}。
\begin{proofsketch}
定义 \(\Sigma X=0\amalg_X0\)、\(\Omega X=0\times_X0\)。同一方块既为拉回又为推出，使 \(\Omega\Sigma X\simeq X\)、\(\Sigma\Omega X\simeq X\)，故移位是等价。映射空间因此可以反复写成循环空间，其连通分支具有交换群结构；有限积与有限余积由同一双笛卡儿方块识别为双积。于是同伦范畴为加性范畴。

对 \(f:X\to Y\) 取余纤维 \(C_f=Y\amalg_X0\)，得到
\(X\to Y\to C_f\to\Sigma X\)。恒等映射的余纤维为零，零映射也有余纤维，因而满足完成三角的公理。将这个推出方块同时看作拉回，可以把余纤维序列改写为纤维序列并旋转；改变循环方向给出旋转中的负号。一个在同伦范畴中交换的方块可以选取实现其交换性的同伦，再用推出的泛性质补出余纤维之间的态射，这给出三角态射公理。

对可复合的 \(X\xrightarrow fY\xrightarrow gZ\)，推出的粘合律给出
\[
 C_f\longrightarrow C_{gf}\longrightarrow C_g
 \longrightarrow\Sigma C_f,
 \qquad
 C_{gf}\amalg_{C_f}0\simeq Z\amalg_Y0=C_g.
\]
把它与 \(f,g,gf\) 各自的余纤维方块组合，就是八面体图。各面同时来自同一推出图，故所需复合相容。这里省略无穷范畴模型中推出粘合与高阶同伦相容性的技术验证；对复形的具体模型，\secref{b4:sec:1602}已经用锥矩阵验证同一构造与负号。这一概要从稳定结构导出三角公理，没有断言一个给定三角范畴会自动具有唯一的稳定增强。
\end{proofsketch}

复形给出最贴近本书的模型。这里固定域 \(k\)，考虑 \(k\) 向量空间的复形。由其 DG 结构构造相应稳定无穷范畴后，映射空间在零映射处的同伦群满足
\[
\pi_i\operatorname{Map}(X,Y)
\simeq H^{-i}\operatorname{Hom}_k^\bullet(X,Y)
\simeq H^{-i}\mathbf R\operatorname{Hom}_k(X,Y),
\qquad i\ge0.
\]
在 \(i=0\) 时，这就是导出范畴中的态射群。第二个同构使用底域条件：由\cref{b4:exm:B-field-complex-model}，零调复形可缩，普通 Hom 已计算导出 Hom。一般环上必须先选择适当的导出 DG 模型，不能直接对普通 Hom 使用同一公式。
\begin{proofsketch}
令 \(H=\operatorname{Hom}_k^\bullet(X,Y)\)。先取非正次的典范截断，再改成非负次链复形：\(A_n=H^{-n}\) 对 \(n>0\)，\(A_0=\ker(H^0\to H^1)\)。将它单纯化时，第 \(m\) 层取
\[
 (\Gamma A)_m=\bigoplus_{[m]\twoheadrightarrow[r]}A_r.
\]
面映射与退化映射由有限有序集映射的满射、单射分解及 \(A\) 的微分给出。除去退化分量以后，正规化复形重新得到 \(A\)；其余分量配成可缩项。所以该单纯交换群的同伦群就是 \(H_i(A)=H^{-i}(H)\)。零分支的描述正是余圈模余边。

DG 神经把这一单纯化与态射复合的高阶相容性组合成映射空间，其取正规化后的上述复形不变，因而给出所列公式。概要省略面映射恒等式和 DG 神经映射空间与此单纯化模型的比较，不省略截断和次数方向；这些是应用公式时必需的条件。底域情形的 Hom 与导出 Hom 已在前面识别。
\end{proofsketch}

例如取 \(X=k\)、\(Y=k[1]\)，总 Hom 仅在 \(-1\) 次为 \(k\)，故映射空间连通而有 \(\pi_1=k\)。普通导出范畴中从 \(k\) 到 \(k[1]\) 只有零态射，却不能据此说相应映射空间没有更高信息。

\ProseParagraphBreak
高阶 \(K\) 理论需要组织大量对象、等价与正合序列之间的相容性，因此稳定无穷范畴为它提供了统一的对象范围。本文止于上述定义与复形例子；完整的增强构造、谱值 \(K\) 理论及其泛性质属于进一步阅读。前面已经证明的 Euler 同构与矩阵 \(K_1\) 计算不依赖这些更高结构。


\begin{chaptersummary}
正合范畴把“哪些列可用来作正合论证”作为额外数据。同一加性范畴可以采用分裂结构或更大的结构，因此零调复形、导出范畴以及 Grothendieck 群都可能随之改变。扩张封闭子范畴提供大量自然例子，而不要求全部核和余核都留在其中。

\ProseParagraphBreak
\(K_0\) 表达可加不变量的泛性质。对有限生成投射模，有界复形的项的交错和在三角关系下不变，给出 \(K_0(R)\simeq K_0\bigl(\operatorname{Perf}(R)\bigr)\)。这种等式保留可加信息，仍可能把非零导出对象送到零类；同调对象不在投射模范畴中时，也不能把它们的类直接写进 \(K_0(R)\)。

\ProseParagraphBreak
\(K_1\) 从稳定自同构群的交换化出发，初等矩阵恰组成其交换子群。域和整数上的计算由行列式及消元完成。高阶 \(K_n\) 需要记录对象、正合序列与复合的相容性，\(Q\) 构造将这些资料组成空间；低阶比较与局部化的概要说明这些空间怎样恢复代数群并给出长正合列，神经实现和群完成的技术论证另见所引著作。

\ProseParagraphBreak
模型结构提供可计算的替换与提升；稳定无穷范畴进一步保留映射空间及高阶同伦。域上的复形使这些思想有具体模型，但一般环上的零调复形未必可缩。三角范畴是这些增强结构的重要结果之一，不能反过来不加条件地恢复全部相干数据。
\end{chaptersummary}

\BookExercises

\begin{exercise}\label{b4:ex:B-exponent-extension}
在有限交换 \(p\)-群范畴中，比较“所有 \(p\)-幂阶群”与“所有被 \(p\) 零化的群”这两个对象族的扩张封闭性。后一族本身是否仍能有正合结构？
\end{exercise}
\begin{solution}
若短正合列两端阶均为 \(p\) 的幂，中项阶为两端阶的乘积，仍为 \(p\) 的幂，所以第一族扩张封闭。后一族不扩张封闭，因为
\(0\rightarrow\mathbb Z/p\mathbb Z\rightarrow\mathbb Z/p^2\mathbb Z\rightarrow\mathbb Z/p\mathbb Z\rightarrow0\)
的两端被 \(p\) 零化，而中项不是。不过被 \(p\) 零化的有限交换群等价于有限维 \(\mathbb F_p\)-向量空间，当然可以配备分裂正合结构；其自身的通常向量空间短正合列也全部分裂。失败的是从较大环境继承全部扩张的封闭性，不是说这些对象无法组成任何正合范畴。
\end{solution}

\begin{exercise}\label{b4:ex:B-nonbalanced-free}
在有限自由 \(\mathbb Z\)-模范畴中证明乘 \(2:\mathbb Z\rightarrow\mathbb Z\) 同时是范畴意义的单态射和满态射，却不是同构。它为什么不是分裂结构中的可容许满态射？
\end{exercise}
\begin{solution}
单态射性来自底层映射单射。若两个到有限自由模 \(F\) 的映射 \(f,g\) 满足 \(2(f-g)=0\)，则 \(F\) 无挠，使 \(f-g=0\)，故乘 \(2\) 也是范畴意义的满态射。它的逆不可能为整数线性映射，所以不是同构。分裂结构的可容许满态射必须有右逆，而乘 \(2\) 没有右逆。此例同时说明该加性范畴不是交换范畴，也说明一般满态射不能替代可容许满态射。
\end{solution}

\begin{exercise}\label{b4:ex:B-pullback-integers}
把短正合列 \(0\rightarrow\mathbb Z\xrightarrow{2}\mathbb Z\rightarrow\mathbb Z/2\mathbb Z\rightarrow0\) 沿约化映射 \(\mathbb Z/4\mathbb Z\rightarrow\mathbb Z/2\mathbb Z\) 拉回。求中项的同构类型，并判断拉回列是否分裂。
\end{exercise}
\begin{solution}
中项为 \(P=\bigl\{(a,\bar b)\in\mathbb Z\times\mathbb Z/4\mathbb Z:a\equiv b\pmod2\bigr\}\)。映射
\[
\mathbb Z\oplus\mathbb Z/2\mathbb Z\longrightarrow P,\qquad
(a,\bar e)\longmapsto(a,\overline{a+2e})
\]
是同构：第二坐标与 \(a\) 同奇偶，故 \(e\) 模 \(2\) 唯一。在这些坐标中，左边的包含为
\(n\mapsto(2n,\overline{-n})\)，右边的满射为
\((a,\bar e)\mapsto\overline{a+2e}\)。
若分裂，\(\mathbb Z/4\mathbb Z\) 会嵌入中项；但中项的挠子群只有 \(\mathbb Z/2\mathbb Z\)，不含阶为 \(4\) 的元素，故不分裂。
\end{solution}

\begin{exercise}\label{b4:ex:B-two-derived-zeroes}
设 \(k\) 为域，\(R=k[\varepsilon]/(\varepsilon^2)\)，把 \(k\) 视为 \(R/(\varepsilon)\)。计算 \(\operatorname{Ext}_R^1(k,k)\)，并证明短正合列
\[
 \xi:\quad 0\longrightarrow k\xrightarrow{\iota}R
 \xrightarrow{\pi}k\longrightarrow0,
 \qquad \iota(c)=c\varepsilon,
\]
给出非零扩张类。把三个非零项视为复形，比较其在通常正合结构和分裂正合结构中的零调性。为什么这个复形在通常导出范畴中为零，却仍有非零的连接态射 \(k\to k[1]\)？
\end{exercise}
\begin{solution}
乘 \(\varepsilon:R\to R\) 的核与像均为 \(\varepsilon R\)，所以
\[
 \cdots\xrightarrow{\varepsilon}R\xrightarrow{\varepsilon}R
 \xrightarrow{\varepsilon}R\xrightarrow{\pi}k\longrightarrow0
\]
是自由消解。施行 \(\operatorname{Hom}_R(-,k)\) 后微分全为零，故 \(\operatorname{Ext}_R^1(k,k)\cong k\)。若 \(\pi\) 有 \(R\)-线性截面 \(s\)，则 \(\varepsilon s(1)=s(\varepsilon\cdot1)=0\)，迫使 \(s(1)\in\varepsilon R\)，但这又使 \(\pi s(1)=0\)，与截面条件矛盾。由\cref{b4:thm:ext1-extensions}，\(\xi\) 的类非零。

三项复形在通常结构中零调，在通常导出范畴中为零。在分裂结构中，若它可缩，末端收缩方程便给出 \(\pi\) 的截面，故它不可缩；由\cref{b4:prop:B-split-acyclic-contractible}，它也不是分裂零调复形，在分裂导出范畴中非零。另一方面，短正合列给出好三角 \(k\to R\to k\xrightarrow{\delta}k[1]\)。由\cref{b4:thm:derived-ext-identification}及其后的扩张对应，\(\delta\) 正是非零扩张类所对应的态射。零调的三项复形与这个好三角记录的数据不同：前者是一个复形对象，后者还保留三项之间的连接态射，前者为零并不迫使后者分裂。
\end{solution}

\begin{exercise}\label{b4:ex:B-finite-p-kzero}
设 \(p\) 为素数，\(M\) 为交换群。分类所有取值于 \(M\)、只依赖有限交换 \(p\)-群的同构类且对直和可加的函数 \(v\)。其中哪些函数还对全部短正合列可加？据此判断“循环直和因子的个数”是否为通常正合结构的可加不变量。
\end{exercise}
\begin{solution}
由\cref{b4:prop:B-p-group-kzero}及 Grothendieck 群的泛性质，可任意指定 \(w_n=v(\mathbb Z/p^n\mathbb Z)\in M\)。若
\[
A\simeq\bigoplus_{n\ge1}(\mathbb Z/p^n\mathbb Z)^{a_n},
\]
则 \(v(A)=\sum_na_nw_n\)，且这些公式给出全部对直和可加的函数。若还对短正合列可加，利用 \(n\ge2\) 时的列
\[
 0\longrightarrow\mathbb Z/p^{n-1}\mathbb Z
 \xrightarrow{\ \bar a\mapsto\overline{pa}\ }
 \mathbb Z/p^n\mathbb Z\longrightarrow\mathbb Z/p\mathbb Z
 \longrightarrow0
\]
得到 \(w_n=w_{n-1}+w_1\)，故 \(w_n=nw_1\)。反之，满足这个条件时 \(v(A)=(\log_p|A|)w_1\)；短正合列中的群阶相乘，因而 \(v\) 可加。取 \(M=\mathbb Z\)、每个 \(w_n=1\)，得到循环因子的个数。但 \(\mathbb Z/p^2\mathbb Z\) 只有一个循环因子，上述 \(n=2\) 的短正合列两端各有一个，故这个函数只对直和可加。
\end{solution}

\begin{exercise}\label{b4:ex:B-weighted-dimension}
设 \(k\) 为域，\(M\) 为交换群。分类所有取值于 \(M\)、对有限维左 \(k\times k\)-模的短正合列可加的不变量。
\end{exercise}
\begin{solution}
模由两个有限维 \(k\)-向量空间 \(V_1,V_2\) 组成，短正合列逐分量正合。由\cref{b4:prop:B-basic-k0}，\(K_0(k\times k)=\mathbb Z^2\)。根据泛性质，每个可加不变量唯一由 \(m_1,m_2\in M\) 确定，公式为
\(v(V_1,V_2)=(\dim V_1)m_1+(\dim V_2)m_2\)。
两个基对象 \((k,0),(0,k)\) 的取值可任意指定，所以没有额外约束。
\end{solution}

\begin{exercise}\label{b4:ex:B-euler-same-different}
对不同素数 \(p,q\)，把 \(\mathbb Z/p\mathbb Z,\mathbb Z/q\mathbb Z\) 放在零次。证明它们在 \(\operatorname{Perf}(\mathbb Z)\) 中不同构，但其 \(K_0\) 类相同。
\end{exercise}
\begin{solution}
两者都由两项自由消解 \(\mathbb Z\xrightarrow{p}\mathbb Z\) 与 \(\mathbb Z\xrightarrow{q}\mathbb Z\) 表示，所以各自的 Euler 类都为 \([\mathbb Z]-[\mathbb Z]=0\)。但导出同构必在零次上同调诱导群同构，\(\mathbb Z/p\mathbb Z\) 与 \(\mathbb Z/q\mathbb Z\) 的阶不同，故不同构。即使局限于完美对象，\(K_0\) 类也不能替代对象本身。
\end{solution}

\begin{exercise}\label{b4:ex:B-nonprojective-cohomology}
设 \(k\) 为域，\(R=k[\varepsilon]/(\varepsilon^2)\)，考虑次数为 \(-1,0\) 的左 \(R\)-模复形 \(R\xrightarrow{\varepsilon}R\)。计算其上同调和在 \(K_0(R)\) 中的 Euler 类，并说明为何不能直接在此群中写上同调模的类。
\end{exercise}
\begin{solution}
两项上同调分别为 \(\varepsilon R\simeq k\) 与 \(R/\varepsilon R\simeq k\)。复形由两个有限自由模组成，Euler 类为零。但 \(k\) 不是投射 \(R\)-模：若商映射 \(R\rightarrow k\) 有 \(R\)-线性截面，\(1\in k\) 的像 \(u\) 必满足 \(\varepsilon u=0\)，故 \(u\in\varepsilon R\)，又使商像为零，与截面矛盾。因此 \([k]\) 不是按有限生成投射模定义的 \(K_0(R)\) 的对象生成元；若要在 \(G_0(R)\) 中使用上同调公式，则应明确更换群及相应的正合范畴。
\end{solution}

\begin{exercise}\label{b4:ex:B-scalar-extension-k0}
设 \(f:R\rightarrow S\) 为含幺环同态。证明 \(P\mapsto S\otimes_RP\) 诱导 \(K_0(R)\rightarrow K_0(S)\)，无需假设 \(S\) 为平坦右 \(R\)-模。计算域的对角嵌入 \(k\rightarrow k\times k\) 及两个投影在 \(K_0\) 上的作用。
\end{exercise}
\begin{solution}
若 \(P\oplus Q\simeq R^n\)，则
\((S\otimes_RP)\oplus(S\otimes_RQ)\simeq S^n\)，所以扩标量保持有限生成投射模。投射模范畴中的可容许短正合列全分裂，任意加性函子均保持它们，故不需要一般短正合列上的平坦性。由泛性质得到群同态。对角嵌入把秩一自由模送到 \((k,k)\)，故在 \(K_0\) 上为 \(n\mapsto(n,n)\)。两个投影分别取第一、第二坐标，其与对角的复合均为恒等。
\end{solution}

\begin{exercise}\label{b4:ex:B-matrix-kzero}
设 \(k\) 为域，\(R=M_n(k)\)。证明 \(K_0(R)\simeq\mathbb Z\)，并求正则左模 \(R\) 在以简单列模 \(k^n\) 为生成元的坐标。
\end{exercise}
\begin{solution}
对任意左 \(R\)-模 \(M\)，矩阵单位给出分解 \(M=\bigoplus_iE_{ii}M\)，而 \(E_{i1},E_{1i}\) 使各分量与 \(E_{11}M\) 同构。因此映射
\(k^n\otimes_k E_{11}M\rightarrow M\)、\(e_i\otimes v\mapsto E_{i1}v\)
为 \(R\)-线性同构，逆由各矩阵单位投影给出。有限生成模对应有限维 \(E_{11}M\)，所以都是有限份简单列模的直和，且都投射；\(K_0(R)\) 因而为 \(\mathbb Z\)。正则模的 \(E_{11}R\) 是任意第一行组成的 \(n\)-维空间，故 \(R\) 是 \(n\) 份列模的直和，其类为 \(n\)，不是这个坐标中的 \(1\)。
\end{solution}

\begin{exercise}\label{b4:ex:B-small-matrix-stabilization}
证明 \(\operatorname{GL}_3(\mathbb F_2)\) 是完美群。再证明稳定化包含 \(j:\operatorname{GL}_2(\mathbb F_2)\to\operatorname{GL}_3(\mathbb F_2)\) 不存在群同态 \(\rho\) 满足 \(\rho j=1\)。
\end{exercise}
\begin{solution}
在 \(\mathbb F_2\) 上，唯一非零标量为 \(1\)，所以消元无须非平凡的行倍乘。行交换也可写成 \(e_{ij}(1)e_{ji}(1)e_{ij}(1)\)，因此所有可逆三阶矩阵均由初等矩阵生成。对不同指标 \(i,j\)，取剩下的指标 \(r\)，由\cref{b4:lem:B-whitehead}中的计算有
\[
 e_{ij}(1)=\bigl[e_{ir}(1),e_{rj}(1)\bigr].
\]
每个生成元都在交换子群中，故 \(\operatorname{GL}_3(\mathbb F_2)\) 完美。若 \(\rho j=1\)，则 \(\rho\) 满射。完美群的满同态像仍完美，因为群同态把交换子送到像中的交换子；但\cref{b4:prop:B-stabilization-example}给出 \(\operatorname{GL}_2(\mathbb F_2)\) 的交换化为 \(C_2\)，它不是完美群，矛盾。
\end{solution}

\begin{exercise}\label{b4:ex:B-integer-matrix-kone}
在 \(K_1(\mathbb Z)\) 中计算矩阵
\[
A=\begin{pmatrix}2&1\\1&1\end{pmatrix},
\qquad B=\begin{pmatrix}-1&0\\0&1\end{pmatrix}
\]
的类，并给出 \(A\) 的初等矩阵分解。
\end{exercise}
\begin{solution}
直接相乘可得 \(A=e_{12}(1)e_{21}(1)\)，因此 \(A\in E(\mathbb Z)\)，其 \(K_1\) 类为单位类。\(B\) 的行列式为 \(-1\)，由\cref{b4:thm:B-k1-field-integers}它给出 \(K_1(\mathbb Z)\) 的非平凡元素。其平方为恒等矩阵，正符合这个群为二阶群的结论。
\end{solution}

\begin{exercise}\label{b4:ex:B-product-kone}
证明含幺环的有限直积满足 \(K_1(R\times S)\simeq K_1(R)\times K_1(S)\)。据此计算 \(K_1(\mathbb F_3\times\mathbb F_5)\)。
\end{exercise}
\begin{solution}
逐项投影给出 \(\operatorname{GL}_n(R\times S)\simeq\operatorname{GL}_n(R)\times\operatorname{GL}_n(S)\)。稳定化后仍是直积，因为一对有限大小矩阵可同时补到同一大小。直积群的交换子逐分量计算，交换子群亦为两者交换子群的直积。取商得到结论。有限域乘法群给出
\(K_1(\mathbb F_3\times\mathbb F_5)\simeq C_2\times C_4\)；例如各因子的生成元可取 \(-1\) 与 \(2\)。
\end{solution}

\begin{exercise}\label{b4:ex:B-polynomial-kone}
把正文对整数的消元论证推广到 \(k[t]\)，证明 \(K_1\bigl(k[t]\bigr)\simeq k^\times\)，其中 \(k\) 为域。
\end{exercise}
\begin{solution}
行列式给出到 \(k[t]^\times=k^\times\) 的满射。若矩阵行列式为 \(1\)，第一列各多项式生成单位理想。用多项式除法及带符号交换执行 Euclid 算法，将第一列化为常数单位与零组成的列，再用 \(\operatorname{diag}(u,u^{-1})\) 将该单位化成 \(1\)。这种二阶对角块由\cref{b4:lem:B-whitehead}中的分解属于初等子群。列加法清掉第一行后，按矩阵大小归纳。故行列式的核恰是初等子群，稳定商为 \(k^\times\)。
\end{solution}

\begin{exercise}\label{b4:ex:B-q-zero-arrows}
令 \(\mathcal E\) 为有限维 \(\mathbb F_q\)-向量空间范畴，\(\theta:B\to C\) 为由 \(B\xleftarrow{p}N\hookrightarrow C\) 表示的 \(Q\)-态射。按\cref{b4:prop:B-q-zero-subspaces}把 \(\operatorname{Hom}_{Q\mathcal E}(0,B)\) 识别为 \(B\) 的子空间集合，计算后复合 \(\theta\) 的作用，证明其为单射，并描述其像。对 \(N=C=\mathbb F_q^3\)、\(B=\mathbb F_q^2\)，令 \(p\) 为前两坐标的投影，分别计算像中各维数子空间的数目。
\end{exercise}
\begin{solution}
子空间 \(U\subseteq B\) 对应 \(0\leftarrow U\hookrightarrow B\)。与 \(\theta\) 复合时，拉回为 \(U\times_BN\)，经 \(N\hookrightarrow C\) 识别为 \(p^{-1}(U)\)。由于 \(p\) 满射，\(p(p^{-1}(U))=U\)，所以这个对应单射。其像恰为
\[
 \{\,W\subseteq C:\ker p\subseteq W\subseteq N\,\}.
\]
事实上，逆像显然满足这两个包含；反之，若 \(\ker p\subseteq W\subseteq N\)，则 \(p^{-1}(p(W))=W\)，因两个有相同商像的向量之差属于 \(\ker p\subseteq W\)。具体投影的核是第三坐标轴 \(L\)，所以像中各子空间的维数为一、二、三，对应 \(B\) 中的零子空间、一维子空间及全空间，数目依次为 \(1,q+1,1\)。这里复合使所有子空间的维数增加一，却保持它们之间的包含关系。
\end{solution}

\begin{exercise}\label{b4:ex:B-q-dimension}
在上一题的范畴中，对 \(a,b\ge0\) 判断从 \(\mathbb F_q^a\) 到 \(\mathbb F_q^b\) 的 \(Q\)-态射何时存在。进一步用 \(b\)-维空间的 \(m\)-维子空间数 \(N_{b,m}\) 写出其态射总数。
\end{exercise}
\begin{solution}
中间子空间 \(M\subset\mathbb F_q^b\) 必有满射到 \(\mathbb F_q^a\)，故 \(a\le\dim M\le b\)。因此存在当且仅当 \(a\le b\)。固定 \(M\) 的维数为 \(m\) 后，满射数等于对偶空间中 \(a\) 个线性无关向量组的数目，即 \(\prod_{j=0}^{a-1}(q^m-q^j)\)。所以总数为
\[
\sum_{m=a}^bN_{b,m}\prod_{j=0}^{a-1}(q^m-q^j).
\]
取中间对象同构类已经把它识别为实际子空间及其指定满射，不须再除以一次 \(\operatorname{GL}_m(\mathbb F_q)\)；重复除法会错误地丢掉态射。
\end{solution}

\begin{exercise}\label{b4:ex:B-localization-arithmetic}
在素数赋值映射 \(\mathbb Q^\times\rightarrow\bigoplus_p\mathbb Z\) 下，求 \(360/49\) 的像。再求局部化环 \(\mathbb Z[1/30]\) 的单位群，并解释为何只有三个素数坐标可以任意变化。
\end{exercise}
\begin{solution}
\(360/49=2^3 3^2 5\,7^{-2}\)，故非零坐标为
\(v_2=3,v_3=2,v_5=1,v_7=-2\)。在 \(\mathbb Z[1/30]\) 中，分母只含 \(2,3,5\)。一个有理数及其倒数同时属于该环，当且仅当它除符号以外也只含这三个素因子。因此单位恰为
\(\{\pm2^a3^b5^c:a,b,c\in\mathbb Z\}\)，与 \(C_2\times\mathbb Z^3\) 同构。其他素数的赋值必须为零，正是倒数仍须在此环中的条件。
\end{solution}

\begin{exercise}\label{b4:ex:B-model-nonfield}
说明若在 \(\operatorname{Ch}(\mathbb Z)\) 上仍以全部逐次单射、逐次满射、拟同构分别作余纤维化、纤维化、弱等价，就不满足模型范畴的提升公理。
\end{exercise}
\begin{solution}
取零调但不可缩的三项复形
\(X=(\mathbb Z\xrightarrow{2}\mathbb Z\rightarrow\mathbb Z/2\mathbb Z)\)。
不可缩是因为末端满射没有整数线性截面。若所给三类满足模型公理，\(0\rightarrow X\) 为平凡余纤维化。\cref{b4:exm:B-field-complex-model}证明中的锥构造只使用复形的微分与恒等映射，因而这里仍适用于整数环；由\cref{b4:prop:cone-sign-compatibility}，\(C(X)=\operatorname{Cone}(1_X)\) 可缩，其移位上的投影给出逐次满射
\(q:C(X)[-1]\rightarrow X\)，所以它是纤维化。以 \(X\rightarrow X\) 为底边恒等映射、顶边为零，提升公理给出截面 \(s:X\rightarrow C(X)[-1]\)，满足 \(qs=1_X\)。可是 \(C(X)[-1]\) 可缩，若 \(h\) 为其收缩，则 \(qhs\) 使 \(X\) 也可缩，矛盾。故这种三类态射不能原样用于整数环。
\end{solution}

\begin{exercise}\label{b4:ex:B-shifted-mapping}
设 \(k\) 为域，\(a,b\in\mathbb Z\)。在域上复形的稳定模型中，计算
\[
\operatorname{Map}\bigl(k[a],k[b]\bigr)
\]
的全部非负次同伦群，并注明移位方向。
\end{exercise}
\begin{solution}
约定 \(X[1]^n=X^{n+1}\)，所以 \(k[a]\) 仅在 \(-a\) 次非零。总 Hom 中，一个 \(r\) 次映射须满足 \(-a+r=-b\)，故总 Hom 仅在 \(r=a-b\) 次为 \(k\)，微分为零。\subsecref{b4:subsec:B04:02}中的识别为 \(\pi_i\operatorname{Map}(X,Y)\cong H^{-i}\operatorname{Hom}_k^\bullet(X,Y)\)，所以非零恰在 \(-i=a-b\)，即 \(i=b-a\ge0\)。因此
\[
\pi_i\operatorname{Map}\bigl(k[a],k[b]\bigr)
=\begin{cases}k,&i=b-a\ge0,\\0,&\text{其余}\;i\ge0.\end{cases}
\]
这里 \(\pi_0\) 按映射空间自然的交换群结构理解。若 \(a=b\)，分支群为 \(k\) 而更高同伦群为零；若 \(b>a\)，空间连通但在正次 \(b-a\) 处有非零同伦；若 \(b<a\)，所有非负次同伦群都为零。
\end{solution}

\begin{exercise}\label{b4:ex:B-ordinary-not-stable}
把\cref{b4:prop:B-ordinary-stability-obstruction}推广到任意非零交换范畴，证明其离散映射空间的无穷范畴不是稳定的。再设 \(k\) 为域，\(X,Y\) 为 \(k\)-向量空间复形，写出零映射 \(X\to Y\) 的同伦余纤维，以及相应余纤维序列中的两条后续箭头。与零映射的普通余核比较。
\end{exercise}
\begin{solution}
在任意非零交换范畴中取 \(V\ne0\)。由 \(V\) 和三个零对象组成的方块仍是推出方块，因为从两个零对象出发的态射唯一；但零对象在零对象上的拉回为零，不是 \(V\)。这与向量空间的证明相同。

零映射的锥在次数 \(n\) 上为 \(Y^n\oplus X^{n+1}\)，微分为
\[
 \mathrm d(y,x)=(\mathrm d_Yy,-\mathrm d_Xx),
\]
所以同伦余纤维为 \(Y\oplus X[1]\)。相应序列为
\[
 X\xrightarrow{0}Y\xrightarrow{\ y\mapsto(y,0)\ }
 Y\oplus X[1]\xrightarrow{\ (y,x)\mapsto x\ }X[1],
\]
后两条箭头是包含与正投影，符合本卷的锥三角约定。普通余核只有 \(Y\)，同伦余纤维还保留 \(X[1]\)。若 \(X\) 的上同调非零，这个附加项在导出范畴中不能删去；若 \(X\) 零调，它才在导出范畴中成为零。
\end{solution}
